\section{Evaluación de políticas y diseño experimental}

\subsection{Métricas fuera de línea y comparación}

La evaluación fuera de línea suele realizarse en un entorno holdout distinto del conjunto de entrenamiento, para asegurar que la política no se haya sobreajustado. Las métricas más comunes incluyen la reward acumulada, la longitud del episode, la estabilidad (desviación estándar) y la sample efficiency. La siguiente tabla muestra qué observa cada una de estas métricas:

\begin{table}[H]
\centering
\begin{tabular}{lp{8cm}}
\toprule
Métrica & Observaciones \\
\midrule
Reward acumulada & Refleja el grado en que la política completa la tarea a partir de un punto de inicio fijo \\
Episode length & Evalúa si aprendió una ruta más concisa / eficiente \\
Sample efficiency & Magnitud de la mejora en reward por unidad de interacción con el entorno \\
Estabilidad & Fluctuación de la reward entre distintas seeds; si es demasiado grande, implica riesgos en el despliegue \\
\bottomrule
\end{tabular}
\caption{Métricas habituales de evaluación fuera de línea}
\end{table}

\begin{importantbox}{Equilibrio entre sample efficiency y reward ceiling}
En ocasiones, el beneficio marginal de aumentar la reward acumulada es muy bajo, pero la sample efficiency todavía no se ha estabilizado, lo que indica que el agent sigue explorando de forma repetida. Se recomienda dar seguimiento a ambas en la curva de iteraciones y, durante el periodo de plateau, introducir una exploration más intensa o una auxiliary task.
\end{importantbox}

\subsection{Comparación en línea e intervención}

En el entorno de producción, una política nueva puede probarse mediante una comparación del tipo champion/challenger o A/B. Debe garantizarse que el challenger opere solo con una fracción pequeña del tráfico y tener preparado un plan de reversión.

\begin{knowledgebox}{A/B vs Champion/ Challenger}
A/B: el tráfico en línea se dirige directamente al challenger según una proporción; es adecuado cuando la diferencia entre la behavior policy y el baseline es pequeña. Champion/Challenger: el challenger se ejecuta en un entorno shadow y, mediante offline evaluation, se decide si se le hace promote.
\end{knowledgebox}

\begin{warningbox}{Desplazamiento de la distribución en la evaluación}
En cuanto existe una diferencia en el ruido de observación entre el entorno de entrenamiento y el de despliegue, la señal de las métricas fuera de línea puede distorsionarse. En sistemas críticos para la seguridad (como la conducción autónoma), debe darse prioridad a construir una biblioteca de scenarios para realizar stress tests.
\end{warningbox}

\subsection{Resumen de la sección}

Esta sección repasó el sistema de evaluación de Q-Learning, subrayó la importancia de atender a la vez la reward acumulada, la eficiencia de muestreo y la estabilidad, y recomendó establecer pruebas comparativas y stress tests antes del despliegue.
